\documentclass[12pt, a4paper]{ctexart}
\usepackage[margin=2cm]{geometry}
\usepackage{libertine}

\usepackage{titlesec}
\usepackage{zhnumber}
\titleformat*{\section}{\Large\bfseries\raggedright}
\renewcommand{\thesection}{\zhnum{section}}
\renewcommand{\thesubsection}{\arabic{subsection}}

\setcounter{secnumdepth}{4}
\renewcommand{\theparagraph}{\thesubsubsection.\arabic{paragraph}}
\renewcommand{\paragraph}[1]{
    \refstepcounter{paragraph}
    {\bfseries\theparagraph\quad#1}\par
    \vspace{2pt}
    \noindent
}

\usepackage{enumitem}
\setlist[enumerate]{itemsep=2pt, parsep=0pt, partopsep=0pt, topsep=2pt}
\setlist[itemize]{itemsep=2pt, parsep=0pt, partopsep=0pt, topsep=2pt}
\linespread{1.2}

\usepackage{graphicx}

\usepackage{minted}
\setminted{
    breaklines=true, escapeinside=||, fontsize=\small, frame=lines,
    mathescape=$$, numbers=none, style=bw, tabsize=4
}

\usepackage{siunitx}

\begin{document}
    
\pagestyle{plain}
\thispagestyle{empty}

\noindent
\begin{tabular*}{\textwidth}{l @{\extracolsep{\fill}} r @{\extracolsep{6pt}} l}
    \LARGE{\textbf{实验报告}} & 计算机网络 & \textit{Computer Networking} \\
\end{tabular*}\\\\
\begin{tabular*}{\textwidth}{l l}
    \textbf{报告标题: } & \textbf{OSPF 路由协议验证} \\
    \textbf{学号: } & 19240212 \\
    \textbf{姓名: } & 华博文 \\
    \textbf{日期: } & 2025 年 12 月 22 日 \\
\end{tabular*}\\
\rule[2ex]{\textwidth}{2pt}

\section{实验目的}

本实验旨在掌握基于 IPMininet 的 OSPF 路由协议配置方法，理解动态路由协议 OSPF 的工作原理（包括网络通告、区域配置等核心机制）；通过 IPMininet 搭建包含 2 台路由器、2 台主机的网络拓扑，完成设备 IP 地址配置、路由器 OSPF 协议部署及全网连通性测试，强化对 OSPF 在自治系统内动态路由自动化管理作用的认知，同时熟悉 IPMininet 仿真工具在网络实验中的应用流程。

\section{实验内容简要描述}

\begin{enumerate}
    \item \textbf{IPMininet 仿真网络启动}：在 Linux 环境中启动 Open vSwitch 服务，进入 IPMininet 示例目录，查看自定义 OSPF 网络拓扑文件并启动仿真网络；
    
    \item \textbf{主机与路由器 IP 配置}：配置主机网卡的 IP 地址及默认网关，配置路由器双网卡的 IP 地址以实现网段互联；
    
    \item \textbf{路由器 OSPF 协议配置}：通过 \mintinline{bash}`telnet` 连接 OSPF 进程，进入配置模式添加目标网段的网络通告（指定区域），完成 OSPF 协议的部署；
    
    \item \textbf{OSPF 协议验证}：测试跨网段主机的 \mintinline{bash}`ping` 连通性，查看路由器路由表及 OSPF 路由信息，验证 OSPF 动态路由的生效状态。
\end{enumerate}

\section{实验步骤与结果分析}

\subsection{IPMininet 仿真网络启动}

右键打开 Linux 终端，执行命令启动 OVS 服务（为IPMininet提供交换功能支持）：

\begin{minted}{bash}
ovs-vswitchd --pidfile --detach
\end{minted}

进入 IPMininet 主目录（\mintinline{text}`/home/headless/ipmininet`），右键打开终端，执行命令启动目标拓扑的仿真网络：

\begin{minted}{bash}
sudo python3 -m ipmininet.examples --topo=simple_ospf4_network
\end{minted}

终端依次显示设备添加日志，最终进入 Mininet CLI 界面，确认网络启动完成。

\begin{figure}[H]
    \centering
    \includegraphics[width=\textwidth]{./img/screenshot_0000.png}
\end{figure}

\subsection{配置主机 IP 地址}

在 Mininet CLI 中打开 h1 终端：

\begin{minted}{bash}
xterm h1
\end{minted}

执行 \mintinline{bash}`ifconfig` 查看默认 IP（与实验规划不符），执行命令配置 IP 为 \mintinline{text}`10.0.0.0/24` 网段：

\begin{minted}{bash}
sudo ifconfig h1-eth0 10.0.0.2/24
\end{minted}

\begin{figure}[H]
    \centering
    \includegraphics[width=\textwidth]{./img/screenshot_0001.png}
\end{figure}

在 Mininet CLI 中打开 h2 终端：

\begin{minted}{bash}
xterm h2
\end{minted}

同理，执行命令配置 IP 为 \mintinline{text}`10.0.2.0/24` 网段：

\begin{minted}{bash}
sudo ifconfig h2-eth0 10.0.2.2/24
\end{minted}

\begin{figure}[H]
    \centering
    \includegraphics[width=\textwidth]{./img/screenshot_0002.png}
\end{figure}

在 Mininet CLI 中打开 r1 终端：

\begin{minted}{bash}
xterm r1
\end{minted}

执行 \mintinline{bash}`ifconfig` 查看网卡信息，配置双网卡 IP：

\begin{minted}{bash}
sudo ifconfig r1-eth0 10.0.0.1/24
sudo ifconfig r1-eth1 10.0.1.1/24
\end{minted}

\begin{figure}[H]
    \centering
    \includegraphics[width=\textwidth]{./img/screenshot_0003.png}
\end{figure}

在 Mininet CLI 中打开 r2 终端，配置双网卡 IP：

\begin{minted}{bash}
sudo ifconfig r2-eth0 10.0.1.2/24
sudo ifconfig r2-eth1 10.0.2.1/24
\end{minted}

\begin{figure}[H]
    \centering
    \includegraphics[width=\textwidth]{./img/screenshot_0004.png}
\end{figure}

在 h1 终端中添加默认网关（指向 r1 接口）：

\begin{minted}{bash}
sudo route add default gw 10.0.0.1
\end{minted}

同理，在 h2 终端中添加默认网关：

\begin{minted}{bash}
sudo route add default gw 10.0.2.1
\end{minted}

执行 \mintinline{bash}`route -n` 确认配置生效：

\begin{figure}[H]
    \centering
    \includegraphics[width=\textwidth]{./img/screenshot_0005.png}
\end{figure}

\subsection{配置路由器 OSPF 协议}

在 Mininet CLI 中查看 r1 路由表：

\begin{minted}{bash}
r1 route
\end{minted}

发现无目标网段路由，需配置 OSPF：

\begin{figure}[H]
    \centering
    \includegraphics[width=\textwidth]{./img/screenshot_0006.png}
\end{figure}

在 r1 终端中 \mintinline{bash}`telnet` 连接 OSPF 进程（密码：\mintinline{text}`zebra`）：

\begin{minted}{bash}
telnet localhost ospfd
\end{minted}

进入 OSPF 配置模式并添加网络通告：

\begin{minted}{bash}
enable
configure terminal
router ospf
network 10.0.0.0/24 area 0.0.0.0
network 10.0.1.0/24 area 0.0.0.0
write
\end{minted}

\begin{figure}[H]
    \centering
    \includegraphics[width=\textwidth]{./img/screenshot_0007.png}
\end{figure}

执行 \mintinline{bash}`show running-config` 确认配置生效：

\begin{figure}[H]
    \centering
    \includegraphics[width=\textwidth]{./img/screenshot_0008.png}
\end{figure}

在 r2 终端中重复上述步骤，添加网络通告：

\begin{minted}{bash}
network 10.0.1.0/24 area 0.0.0.0
network 10.0.2.0/24 area 0.0.0.0
write
\end{minted}

\begin{figure}[H]
    \centering
    \includegraphics[width=\textwidth]{./img/screenshot_0009.png}
\end{figure}

\subsection{验证 OSPF 协议}

在 h1 终端中 \mintinline{bash}`ping` h2：

\begin{minted}{bash}
ping 10.0.2.2
\end{minted}

在 Mininet CLI 中查看 r1、r2 路由表：

\begin{minted}{bash}
r1 route
r2 route
\end{minted}

可见 OSPF 学习到的动态路由（标志 UG）：

\begin{figure}[H]
    \centering
    \includegraphics[width=\textwidth]{./img/screenshot_0010.png}
\end{figure}

在 r1 的 OSPF 进程中查看 OSPF 路由表：

\begin{minted}{bash}
show ip ospf route
\end{minted}

可看到通过 OSPF 获取的目标网段路由：

\begin{figure}[H]
    \centering
    \includegraphics[width=\textwidth]{./img/screenshot_0011.png}
\end{figure}

\section{实验中遇到的问题及体会}

在启动 IPMininet 时，因未切换至正确目录执行命令，终端提示 ``模块找不到'' 错误。核对实验步骤后切换至 IPMininet 根目录，重新执行命令成功启动网络 —— 这让我认识到网络仿真工具对执行环境的依赖较强，操作前需严格确认路径正确性。

配置 OSPF 时，最初遗漏 \mintinline{bash}`write` 命令导致配置丢失，路由器路由表未更新。重新进入 OSPF 进程并执行保存后，路由表成功生成 —— 这体现了网络配置 ``修改 —— 验证 —— 保存'' 的规范流程，临时配置需固化才能长期生效。

通过本次实验，我直观体会到 OSPF 的动态路由优势：无需手动配置网段路由，仅需通告网络即可自动学习拓扑，大幅降低多网段网络的维护成本。同时，IPMininet 通过代码定义拓扑、结合 OVS 实现交换的方式，让我了解到虚拟网络仿真的灵活架构。

实验也深化了 ``网络配置协同性'' 的认知：从主机 IP、路由器接口到 OSPF 通告，每一步需严格匹配网段规划，任何一处错误都会导致路由失效。这种 ``端到端协同'' 的逻辑，是网络通信正常运行的核心准则。

最后，IPMininet 的虚拟环境为实验提供了低成本实践载体，可快速搭建拓扑、反复调试配置。这一过程不仅提升了我的 Linux 命令与网络配置熟练度，更培养了从路由表、协议报文倒推故障的排查思维，为后续高级路由技术的学习奠定了基础。

\end{document}